\renewcommand\LTcaptype{quadro}
\begin{scriptsize}
\begin{longtable}[c]{|p{5cm}|l|l|l|l|l|l|l|}

\caption{Plano de tarefas do TCC 2}
\addcontentsline{loq}{quadro}{\protect\numberline{\thequadro}Plano de tarefas do TCC 2}
\label{tab:TCC2}\\
\hline
\multicolumn{1}{|c|}{Tarefas} &
  \multicolumn{1}{c|}{Out/2025} &
  \multicolumn{1}{c|}{Nov/2025} &
  \multicolumn{1}{c|}{Dez/2025} &
  \multicolumn{1}{c|}{Fev/2026} &
  \multicolumn{1}{c|}{Mar/2026} &
  \multicolumn{1}{c|}{Abr/2026} &
  \multicolumn{1}{c|}{Mai/2026} \\ \hline
\endhead
%
Aprofundar o referencial teórico &
  X & X &  &  &  &  &  \\ \hline
Elaborar o protótipo do sistema/software &
  & X & X &  &  &  &  \\ \hline
Realizar testes experimentais com cadeirantes &
  &  & X & X &  &  &  \\ \hline
Coletar e analisar os dados obtidos &
  &  &  & X & X &  &  \\ \hline
Ajustar a solução proposta conforme os resultados &
  &  &  &  & X & X &  \\ \hline
Redigir capítulos finais (Metodologia aplicada, Resultados e Discussão, Conclusão) &
  &  &  &  &  & X &  \\ \hline
Preparar a versão final do trabalho &
  &  &  &  &  & X &  \\ \hline
Apresentação do trabalho para a banca examinadora &
  &  &  &  &  &  & X \\ \hline
Limite para a entrega dos ajustes após avaliação da banca examinadora &
  &  &  &  &  &  & X \\ \hline
\caption*{Fonte: Elaborado pelo autor.}
\end{longtable}
\end{scriptsize}
\renewcommand\LTcaptype{quadro}
